\begin{table}
\centering
 \begin{tabular}{| c | c | c | c |}
 \hline
 Parameter(s)                                               		   		& Prior		& Range    \\
\hline
$\sin^{2}\theta_{23}$                                       				&  uniform 	& $[0.3; 0.7]$      \\
$\sin^{2} 2 \theta_{13}$ ($\sin^{2} \theta_{13}$) reactors		&  gauss 		& $0.085 \pm 0.005$ 	\\
$\sin^{2} 2 \theta_{13}$ ($\sin^{2} \theta_{13}$) T2K only	   	&  uniform		& $[0; 0.4]$ 	\\
$\sin^{2} 2 \theta_{12}$  								&  gauss 		& $0.846 \pm 0.021$ \\
$|\Delta m^{2}_{32}|$ (NH) / $|\Delta m^{2}_{13}|$ (IH)      		&  uniform		& $ [2; 3] \times 10^{-3}~\mbox{eV}^2/\mbox{c}^4$		 \\
$\Delta m^{2}_{21}$      								&  gauss 		& $(7.53 \pm 0.18) \times 10^{-5}~\mbox{eV}^2/\mbox{c}^4$		\\
$\delta_{CP}$                                               				&  uniform 	& $[-\pi; +\pi]$ \\
Mass Hierarchy										&  fixed 		& NH or IH \\  
\hline
\end{tabular}
\caption{
Treatment of the oscillation parameters in the 
$\nu/\bar{\nu}$ joint analysis.
All the gaussian priors are from \cite{pdg_2014}.
The parameter $\sin^{2} 2 \theta_{13}$ can be constrained with
the measurement of reactor experiments using the gaussian prior,
otherwise a uniform prior is used
and the measurement relies only on the T2K data.
The nominal values of 
$\sin^2 2 \theta_{13}$, 
$\sin^2 2 \theta_{12}$
and 
$\Delta m^2_{21}$ 
are from \cite{pdg_2014},
while all the other oscillation parameter values corresponds to the most probable values 
obtained by the Bayesian analysis on the T2K run 1-4 neutrino mode data \cite{T2KLongPaper}.
The mass hierarchy is not marginalised but fixed to either NH or IH.
}
\label{tab:osc_params_joint_priors}
\end{table}

\begin{table}
\centering
 \begin{tabular}{| c | c | c |}
 \hline
 Parameter(s)                                               		   		&      Asimov A 	  						 & Asimov B   \\
\hline
$\sin^{2}\theta_{23}$                                       				&      0.528								&  0.45   \\
$\sin^{2} 2 \theta_{13}$ ($\sin^{2} \theta_{13}$) reactors		&      0.085 (0.0217)							&  0.085 (0.0217)	 	\\
$\sin^{2} 2 \theta_{12}$  								&      0.846              							&  0.846 		 \\
$|\Delta m^{2}_{32}|$ (NH) / $|\Delta m^{2}_{31}|$ (IH)      		&     	$2.509 \times 10^{-3}~\mbox{eV}^2/\mbox{c}^4$	&  $2.509 \times 10^{-3}~\mbox{eV}^2/\mbox{c}^4$		 \\
$\Delta m^{2}_{21}$      								&       $7.53 \times 10^{-5}~\mbox{eV}^2/\mbox{c}^4$ 	&  $7.53 \times 10^{-5}~\mbox{eV}^2/\mbox{c}^4$ 			\\
$\delta_{CP}$                                               				&       -1.601								& 0    \\
Mass Hierarchy										&        Normal 	         						&  Normal	 \\  
\hline
\end{tabular}
\caption{
Values of oscillation parameters used to compute the event rates,
systematic effects and sensitivity studies.
Each set of oscillation parameters correspond to a different 
Asimov data set, which is the MC expected distribution in a certain
oscillation hypothesis.
In the Asimov data set A the nominal values of 
$\sin^2 2 \theta_{13}$, 
$\sin^2 2 \theta_{12}$
and 
$\Delta m^2_{21}$ 
are from \cite{pdg_2014},
while all the other oscillation parameter values corresponds to the most probable values 
obtained by the Bayesian analysis on the T2K run 1-4 neutrino mode data \cite{T2KLongPaper}.
The mass hierarchy is not marginalised but fixed to either NH or IH.
The only difference between the Asimov data sets A and B is 
given by $\delta_{CP}$, set to 0, 
and $\sin^{2}\theta_{23}$ 
which corresponds to an arbitrary point in the opposite octant.
}
\label{tab:osc_params_joint_asimov}
\end{table}

